\begin{table}[!ht]
\centering
\caption{Escolha de hiperparâmetros e validação cruzada do XGBoost
\cite{alencar2026a, alencar2026b, alencar2026c, geron2021}. Painel~A: seis configurações comparadas numa partição de validação
\emph{dentro} do treino --- o conjunto de teste permanece intocado. Painel~B: validação
cruzada $k$-\emph{fold} ($k=5$) no treino completo, com média e
desvio-padrão entre \emph{folds}. População completa: 6.159.418 observações de
treino e 1.539.855 de teste; nenhuma amostragem.}
\label{tab:ml_cv}
\small
\begin{tabular}{lccccc}
\multicolumn{6}{l}{\textbf{A. Busca na partição de validação}} \\
\toprule
Profundidade & Taxa de aprendizado & Árvores & $R^2$ & MAE & Sobreajuste \\
\midrule
\textbf{10} & 0,05 & 300 & 0,6427 & 0,3977 & 0,0120 \\
8 & 0,05 & 600 & 0,6395 & 0,3995 & 0,0075 \\
8 & 0,05 & 300 & 0,6352 & 0,4022 & 0,0046 \\
6 & 0,10 & 300 & 0,6322 & 0,4041 & 0,0027 \\
6 & 0,05 & 300 (referência) & 0,6286 & 0,4062 & 0,0019 \\
4 & 0,05 & 300 & 0,6182 & 0,4128 & 0,0011 \\
\bottomrule
\end{tabular}

\vspace{0.6em}
\begin{tabular}{lccc}
\multicolumn{4}{l}{\textbf{B. Validação cruzada no treino (5 \emph{folds})}} \\
\toprule
Configuração & $R^2$ (média $\pm$ dp) & MAE & RMSE \\
\midrule
Escolhida (profundidade 10) & 0,6435 $\pm$ 0,0006 & 0,3978 $\pm$ 0,0002 & 0,5360 $\pm$ 0,0002 \\
Referência (profundidade 6) & 0,6294 $\pm$ 0,0006 & 0,4064 $\pm$ 0,0002 & 0,5466 $\pm$ 0,0002 \\
\bottomrule
\end{tabular}
\normalsize
\par\smallskip
\footnotesize\emph{Como ler:} ``Sobreajuste'' é a diferença entre o $R^2$ de treino e o de
validação --- valores próximos de zero indicam que o modelo não decorou os dados. A
configuração escolhida ganha 1,4 ponto percentual de $R^2$ ao custo de um
sobreajuste maior, mas ainda pequeno (Painel~A). No conjunto de teste (intocado durante a
escolha), $R^2 = 0,6444$; a estabilidade entre \emph{folds} e o erro em reais são
discutidos no texto.
\end{table}
